% SOURCE_AUTHORITY: sga5_fr_workpass.tex, lines 14605--14623 (LF-normalized unit).
% SOURCE_SHA256: 791F4EFFC5E02832D5D77ED03518C8156D6F07E4C8238B03545DB93D883FBB28
\begin{statement}{Proposição 1}
\begin{enumerate}[label=\alph*)]
\item O isomorfismo $\Fr^*_{F/X}$ depende funtorialmente do feixe $F$; em outras palavras, tem-se um isomorfismo funtorial
\[
\Fr^*_{/X}:\fr_X^*\longrightarrow \id_{\widetilde X_{\et}},
\]
onde $\widetilde X_{\et}$ é a categoria dos feixes de conjuntos sobre $X_{\et}$.
\item A correspondência de Frobenius $(\fr_X,\Fr^*_{F/X})$ sobre $(X,F)$ é compatível com a mudança de base em $X$. Ou seja, se $\mathcal F$ denota a categoria fibrada sobre $(\Sch)_p$ dos feixes, para a topologia étale, sobre pré-esquemas variáveis de característica $p$, os funtores $\fr_X^*$, $X\in\Ob(\Sch)_p$, definem um $(\Sch)_p$-funtor cartesiano
\[
\fr^*: \mathcal F\longrightarrow\mathcal F,
\]
e os morfismos $\Fr^*_{F/X}$, $X\in\Ob(\Sch)_p$, $F\in\Ob\widetilde X_{\et}$, definem um $(\Sch)_p$-morfismo funtorial
\[
\Fr^*_{/}:\fr^*\longrightarrow\id_{\mathcal F}
\]
que é um isomorfismo.
\item Reciprocamente, estas propriedades caracterizam $\Fr^*_{/}$, ou seja, existe um único $(\Sch)_p$-morfismo funtorial de $\fr^*$ a $\id_{\mathcal F}$.
\end{enumerate}
\end{statement}
